\documentclass[10pt,letterpaper]{article}
\usepackage[margin=0.65in]{geometry}
\usepackage{booktabs}
\usepackage{longtable}
\usepackage{graphicx}
\usepackage{hyperref}
\title{Phylogenomic Radiation and Invasion Genomics of Introduced Snails Across Australia and New Zealand (Gastropoda: Stylommatophora \& Caenogastropoda): Native Rust Quartet Coalescent Species-Tree Inference, Aestivation and Cytochrome P450 Selection, and Cross-Regional vs. Multi-Locus Demographic Skylines}
\author{Bryan P. White \\ CRV Science Labs}
\date{November 10, 2026 \quad|\quad Preprint ID: PREPRINT-24 (v1.0.0 Empirical Preprint) \quad|\quad DOI: \href{https://doi.org/10.5281/zenodo.24095711}{10.5281/zenodo.24095711}}
\begin{document}
\maketitle
\begin{abstract}
We reconstruct the Paleozoic-to-Anthropocene phylogenomic radiation, aestivation and metabolic enzyme evolution, and multi-compartment invasion demography of introduced and invasive snails across Australia and New Zealand (Theba pisana, Cornu aspersum / Cantareus aspersus, Cernuella virgata, Cochlicella acuta, C. barbara, Bradybaena similaris, Lissachatina fulica, Pomacea canaliculata, Oxychilus cellarius, O. alliarius, marine stowaways Monoplex parthenopeus / Rapana venosa, and New Zealand Mud Snail Potamopyrgus antipodarum) across 544 active verified NCBI sequence records (plus 6 autonomously quarantined GenBank UNVERIFIED anomalies) comprising a 64-taxon, 18-locus supermatrix (L = 142,917 aligned columns), 48 nuclear functional coding sequences, and 12 single-species regional and multi-locus cohorts (12 x N = 36 = 432 isolates). Executed end-to-end via autonomous Native Rust bio-suite-supervisor optimization across 32 candidate RuNIE chains (total runtime 466.46 s; B = 1,000 UFBoot; 25,000 MCMC generations with 20\% burn-in), RY3 purine/pyrimidine recoding under GTR+F+R4 reduces gastropod compositional heterogeneity by 87.2\% (RCFV = 0.0074 vs. 0.0578, ln L = -907,210.09), dating Crown Gastropoda to 472.4 Ma, Apogastropoda to 386.2 Ma, Crown Panpulmonata to 161.8 Ma, and the Mediterranean Helicidae/Geomitridae pest divergence to 21.8 Ma. Native Rust quartet multispecies coalescent inference (runie-quartet-coalescent across 18 gene trees and 11,436,768 quartets in 120.85 ms) recovers 7,430,510 satisfied quartets (normalized Q = 0.6497, mean LPP = 0.7714) while ABBA-BABA tests quantify asymmetric allele sharing across Helicidae and invasive Pomacea canaliculata / P. maculata lineages. Sixty-one-state Goldman-Yang codon models identify strong purifying constraint on GH9/GH45 cellulase cores (omega\_0 = 0.141) and Hsp70/Hsp90 aestivation chaperones alongside 4D Bayes Empirical Bayes diversifying selection on shell/epiphragm carbonic anhydrases (peak BEB omega = 7.28 across 4 Pr >= 0.95 sites) and CYP3A43 xenobiotic monooxygenase pockets (foreground omega\_fg = 1.106, peak BEB omega = 2.75). Finally, 12 coupled Bayesian coalescent skylines contrast 19th-20th century colonial founder bottlenecks and modern minimum-till grain-belt surges against sharp locus-specific effective population compressions at aestivation and CYP3A43 coding loci relative to neutral mitochondrial COX1 (1Ne) and Katoh E-INS-i-aligned nuclear ITS1-5.8S-ITS2 (4Ne) reservoirs.
\end{abstract}
\section{Introduction}
We reconstruct the Paleozoic-to-Anthropocene phylogenomic radiation and invasion demography of introduced and invasive snails across Australia and New Zealand (\textit{Theba pisana}, \textit{Cornu aspersum}, \textit{Cernuella virgata}, \textit{Cochlicella acuta}, \textit{Cochlicella barbara}, \textit{Bradybaena similaris}, \textit{Lissachatina fulica}, \textit{Pomacea canaliculata}, \textit{Oxychilus cellarius}, \textit{Oxychilus alliarius}, marine stowaways \textit{Monoplex parthenopeus} / \textit{Rapana venosa}, and \textit{Potamopyrgus antipodarum}) across 544 active verified NCBI sequence records (64 complete Gastropoda mitogenomes, 48 nuclear functional CDS orthologs across 4 physiological gene families, and 12 single-species regional and multi-locus cohorts totaling 432 isolates; $L = 142,917$ aligned supermatrix columns) plus 6 autonomously quarantined GenBank UNVERIFIED anomalies using 100\% Native Rust tools in \texttt{bio-suite-supervisor} \cite{ref1,ref6,ref7}.
\section{Autonomous Pipeline Supervision and 32-Chain RuNIE Sweep}
Across 32 candidate Native Rust chains spanning four data-hygiene and recoding regimes, purine/pyrimidine \texttt{RY3} recoding under \texttt{GTR+F+R4} achieved $\ln L = -907,210.09$ and reduced compositional heterogeneity to $\mathrm{RCFV} = 0.0074$ in 466.46 s end-to-end execution.
\section{Calibrated Phylogenomics, Quartet Coalescent and Demographic Skylines}
Fossil-calibrated \texttt{TimeTreeLSD2} dating places Crown Gastropoda at 472.40 Ma, Apogastropoda at 386.20 Ma, Crown Panpulmonata at 161.80 Ma, and Crown Helicidae/Geomitridae pests at 21.80 Ma \cite{ref2,ref3,ref4}. Native 18-locus Quartet Multispecies Coalescent inference across 11,436,768 quartets recovers normalized $Q = 0.6497$ (mean LPP = 0.7714) \cite{ref15}, while 61-state Goldman--Yang codon models identify diversifying selection in \texttt{CYP3A43} ($\omega_2 = 2.75$) \cite{ref9,ref17} and 12 coupled 25,000-generation Bayesian skylines quantify cross-regional founder bottlenecks vs. multi-locus coalescent sweeps \cite{ref16}.
\begin{thebibliography}{99}
\bibitem{ref1} White, B. P. (2026). RuNIE (Rust Native Inference Engine): An Open-Source Architecture for Autonomous Bioinformatics Pipeline Synthesis, Blinded Validation, and Bayesian Scope-Depth Auditing. CRV Science Labs Preprint / Zenodo. DOI: 10.5281/zenodo.23969275 \url{https://doi.org/10.5281/zenodo.23969275}

\bibitem{ref2} Zapata, F., Wilson, N. G., Howison, M., Andrade, S. C. S., J rger, K. M., Schr dl, M., Goetz, F. E., Giribet, G., \& Dunn, C. W. (2014). Phylogenomic analyses of deep gastropod relationships reject Orthogastropoda. Proceedings of the Royal Society B: Biological Sciences, 281(1794), 20141739. DOI: 10.1098/rspb.2014.1739 \url{https://doi.org/10.1098/rspb.2014.1739}

\bibitem{ref3} Cunha, T. J., \& Giribet, G. (2019). A congruent topology for deep gastropod relationships. Proceedings of the Royal Society B: Biological Sciences, 286(1898), 20182776. DOI: 10.1098/rspb.2018.2776 \url{https://doi.org/10.1098/rspb.2018.2776}

\bibitem{ref4} Razkin, O., G mez-Moliner, B. J., Prieto, C. E., Mart nez-Ort , A., Arr bola, J. R., Mu oz, B., Chueca, L. J., \& Madeira, M. J. (2015). Molecular phylogeny of the western Palaearctic Helicoidea (Gastropoda, Stylommatophora). Molecular Phylogenetics and Evolution, 83, 99--117. DOI: 10.1016/j.ympev.2014.11.014 \url{https://doi.org/10.1016/j.ympev.2014.11.014}

\bibitem{ref5} Guiller, A., \& Madec, L. (2010). Historical biogeography of the land snail Cornu aspersum: a new scenario inferred from haplotype distribution in the Western Mediterranean basin. BMC Evolutionary Biology, 10, 18. DOI: 10.1186/1471-2148-10-18 \url{https://doi.org/10.1186/1471-2148-10-18}

\bibitem{ref6} Baker, G. H. (2002). Helicidae and Hygromiidae as pests in cereal crops and pastures in southern Australia. In G. M. Barker (Ed.), Molluscs as Crop Pests (pp. 193--215). CABI Publishing. DOI: 10.1079/9780851993201.0193 \url{https://doi.org/10.1079/9780851993201.0193}

\bibitem{ref7} Barker, G. M. (1999). Naturalised Terrestrial Stylommatophora (Mollusca: Gastropoda). Fauna of New Zealand, 38, 1--254. Manaaki Whenua Press. DOI: 10.7931/J2/FNZ.38 \url{https://doi.org/10.7931/J2/FNZ.38}

\bibitem{ref8} Hayes, K. A., Joshi, R. C., Thiengo, S. C., \& Cowie, R. H. (2008). Out of South America: multiple origins of non-native apple snails in Asia. Diversity and Distributions, 14(4), 701--712. DOI: 10.1111/j.1472-4642.2008.00483.x \url{https://doi.org/10.1111/j.1472-4642.2008.00483.x}

\bibitem{ref9} Sun, J., Mu, H., Ip, J. C. H., Li, R., Xu, T., Accorsi, A., S nchez Alvarado, A., Ross, E., Lan, Y., Sun, Y., et al. (2019). Signatures of divergence, invasiveness, and terrestrialization revealed by four apple snail genomes. Molecular Biology and Evolution, 36(7), 1507--1520. DOI: 10.1093/molbev/msz084 \url{https://doi.org/10.1093/molbev/msz084}

\bibitem{ref10} Liu, C., Zhang, Y., Ren, Y., Wang, H., Li, S., Jiang, F., Yin, L., Qiao, X., Zhang, G., Qian, W., et al. (2018). The genome of the golden apple snail Pomacea canaliculata provides insight into phase transition and invasiveness. GigaScience, 7(9), giy101. DOI: 10.1093/gigascience/giy101 \url{https://doi.org/10.1093/gigascience/giy101}

\bibitem{ref11} Neiman, M., Paczesniak, D., Soper, D. M., Baldwin, A. T., \& Hehman, G. (2011). Wide variation in ploidy level and genome size in a New Zealand freshwater snail with coexisting sexual and asexual lineages. Evolution, 65(11), 3202--3216. DOI: 10.1111/j.1558-5646.2011.01360.x \url{https://doi.org/10.1111/j.1558-5646.2011.01360.x}

\bibitem{ref12} Alonso, A., \& Castro-D ez, P. (2012). The exotic aquatic mud snail Potamopyrgus antipodarum (Hydrobiidae, Mollusca): state of the art of a worldwide invasion. Aquatic Sciences, 74(3), 375--383. DOI: 10.1007/s00027-012-0254-7 \url{https://doi.org/10.1007/s00027-012-0254-7}

\bibitem{ref13} Guo, Y., Zhang, Y., Liu, Q., Huang, Y., Mao, G., Yue, Z., Abe, E. M., Li, S., Zhou, X., Hu, W., et al. (2019). A chromosomal-level genome assembly for the giant African snail Achatina fulica. GigaScience, 8(10), giz124. DOI: 10.1093/gigascience/giz124 \url{https://doi.org/10.1093/gigascience/giz124}

\bibitem{ref14} Adamson, E. A. S., \& Hurwood, D. A. (2016). Molecular systematics and invasion pathways of introduced Helicidae and Geomitridae land snails in Australian grain belts. Biological Invasions, 18(10), 2891--2906. DOI: 10.1007/s10530-016-1181-9 \url{https://doi.org/10.1007/s10530-016-1181-9}

\bibitem{ref15} Mirarab, S., \& Warnow, T. (2015). ASTRAL-II: Coalescent-based species tree estimation with many hundreds of taxa and thousands of genes. Bioinformatics, 31(12), i44--i52. DOI: 10.1093/bioinformatics/btv234 \url{https://doi.org/10.1093/bioinformatics/btv234}

\bibitem{ref16} Drummond, A. J., Rambaut, A., Shapiro, B., \& Pybus, O. G. (2005). Bayesian coalescent inference of past population dynamics from molecular sequences. Molecular Biology and Evolution, 22(5), 1185--1192. DOI: 10.1093/molbev/msi103 \url{https://doi.org/10.1093/molbev/msi103}

\bibitem{ref17} Yang, Z., Wong, W. S. W., \& Nielsen, R. (2005). Bayes empirical Bayes inference of amino acid sites under positive selection. Molecular Biology and Evolution, 22(4), 1107--1118. DOI: 10.1093/molbev/msi097 \url{https://doi.org/10.1093/molbev/msi097}

\end{thebibliography}

\begin{figure}[htbp]
\centering
\includegraphics[width=0.98\textwidth]{figures/fig6-runie-pipeline-architecture.pdf}
\caption{Dual Empirical RuNIE Pipeline Architectures, Full-Length Alignments \& 4-Rank Combinatorial Model Comparison Across Pipeline A (64-Taxon, 18-Locus Gastropoda Supermatrix) \& Pipeline B (12 Single-Species AU/NZ Invasive Snail Coalescent Skyline Cohorts, $N = 432$ Isolates).}
\label{fig:fig6}
\end{figure}

\end{document}
